\documentclass[12pt,draftclsnofoot,onecolumn]{IEEEtran}

\usepackage[letterpaper,margin=1in]{geometry}
\usepackage[T1]{fontenc}
\usepackage{newtxtext,newtxmath}
\usepackage{microtype}
\usepackage{cite}
\usepackage{booktabs}
\usepackage{array}
\usepackage{tabularx}
\usepackage{url}
\usepackage[hidelinks]{hyperref}

\setlength{\parskip}{0.35em}
\setlength{\parindent}{1em}
\renewcommand{\arraystretch}{1.0}

\title{Empirical Study on the Characteristics and Evolution of AI Usage in GitHub Repositories: Evidence from Code Comments}
\author{Abdullah Al Mujahid\\\normalsize CS 5001\\May 12, 2026}

\begin{document}

\maketitle

\begin{abstract}
Developers are increasingly incorporating AI tools and large language models (LLMs), such as ChatGPT, GitHub Copilot, Claude, Gemini, and Llama, into everyday software development workflows. While many studies evaluate LLM outputs in controlled settings, less is known about how developers integrate, refine, and maintain AI-assisted code in real projects. This study examines GitHub code comments that explicitly reference AI usage and connects those comments to associated code blocks and later commits. The dataset contains 35,361 AI-referenced comments and 35,278 associated code blocks from Python and JavaScript repositories. A mixed-method approach is used, combining manual open coding, LLM-based annotation, Dawid-Skene expectation-maximization aggregation, topic modeling, and longitudinal analysis. The results show that developers most frequently use AI for code implementation, followed by code enhancement, testing, documentation, and bug fixing. Subsequent commits reveal that AI-assisted code often undergoes feature extension, refactoring, cleanup, corrective fixes, documentation, and testing. Longitudinal analysis further shows a shift from direct implementation toward enhancement, testing, and knowledge support. These findings suggest that AI is becoming an active collaborator in software development, but human oversight remains central to successful integration and maintenance.
\end{abstract}

\section{Introduction}

AI tools and LLMs have become increasingly important in modern software development. Developers use systems such as ChatGPT, Copilot, Claude, Gemini, and Llama-based tools to generate code, explain APIs, improve implementations, write tests, support documentation, and explore design alternatives \cite{barke2023grounded,murali2024ai}. This change is reshaping programming from an exclusively human activity into a collaborative process involving both developers and AI systems.

Most prior studies on AI-assisted programming focus on benchmark evaluations, controlled experiments, or short-term user studies \cite{peng2023impact,du2024evaluating}. These studies are useful for measuring model performance, developer productivity, usability, and code quality. However, they provide limited insight into how developers actually use AI-generated or AI-suggested code after it enters a real project. In practice, code must be integrated with existing modules, adjusted to project conventions, tested, reviewed, debugged, and maintained over time. Therefore, evaluating only the initial AI output gives an incomplete view of AI's practical role in software engineering.

This study addresses that gap by examining in-situ artifacts from public GitHub repositories. Code comments are especially valuable for this purpose because they are located beside the relevant source code and often capture developer reasoning, attribution, uncertainty, or justification \cite{rani2023decade,tan2015code}. Comments such as ``generated by ChatGPT,'' ``suggested by Copilot,'' or ``written with Claude'' provide direct evidence of AI involvement in development work. By linking these comments to nearby code blocks and later commits, this study observes not only how AI-assisted code is introduced, but also how it is modified after integration.

The goal of this study is to characterize how developers use AI in real-world software development, what forms of contribution AI provides, how AI-assisted code changes after introduction, and how AI usage behavior evolves over time. The study is guided by three research questions:

\begin{itemize}
    \item \textbf{RQ1:} How do developers integrate AI into real-world software development workflows, and what forms of contribution does AI make?
    \item \textbf{RQ2:} How do developers subsequently adjust, refine, or extend AI-assisted code after its introduction into projects?
    \item \textbf{RQ3:} How has developers' AI usage behavior evolved over time?
\end{itemize}

\section{Background and Related Work}

This study is grounded in two areas of software engineering research: mining repository artifacts to infer developer intent and studying the use of LLMs in software development. Code comments, commit messages, pull requests, and issue discussions have long been used to understand maintenance behavior, design rationale, and developer activity \cite{hindle2009automated}. Comments can reveal assumptions, implementation intent, limitations, and uncertainty, while commit messages often describe corrective, adaptive, perfective, or coordination-related changes.

Recent research on AI-assisted software engineering has examined code generation, code repair, documentation, testing, program comprehension, and developer productivity \cite{fan2023large}. Tools such as GitHub Copilot and ChatGPT have been studied through benchmark tasks, user studies, and mining-based investigations. Prior work shows that developers use LLMs for coding guidance, debugging, API explanation, framework clarification, documentation, testing, and problem solving \cite{hao2024sharedconversations}. Mining-based studies have also examined AI-generated code in open-source repositories and GitHub collaboration artifacts \cite{grewal2024analyzing}.

This study extends prior work by focusing on AI-referenced code comments and associated code blocks. This perspective provides a finer-grained view of AI use than repository-level metadata or broad project discussions alone. It also enables analysis of post-integration changes through first-change commits. As a result, the study captures AI-assisted development as a lifecycle activity rather than as a single code-generation event.

\section{Methodology}

This study adopts a multi-stage, mixed-method research design that combines large-scale data mining, human qualitative coding, LLM-assisted annotation, probabilistic aggregation, topic modeling, and longitudinal analysis. The analysis focuses on Python and JavaScript because they are widely used in open-source development and strongly represented in modern AI-related programming activity \cite{so2024tech}.

The data collection process began with GitHub Code Search API queries designed to identify source code comments that explicitly referenced AI usage. Query templates were constructed by combining AI-related keywords, including ChatGPT, GPT, Copilot, Claude, Gemini, and Llama, with generative action verbs such as generated, suggested, created, written, authored, and assisted. Connector terms such as by, from, with, and using were also included to capture phrases such as ``generated by ChatGPT'' and ``written with Copilot.'' This produced 480 search query patterns.

The initial search returned 66,960 matches across Python and JavaScript files. These results were filtered to retain only cases where the AI-related phrase appeared inside an actual comment span rather than in code, string literals, identifiers, or unrelated text. File histories were then traced to identify the earliest commit in which each matched comment appeared. Only comments introduced between December 2022 and March 2026 were retained. After filtering, the dataset contained 35,361 comments from 26,563 source files in 12,962 repositories.

Associated code blocks were extracted from the introductory commit patches. For each comment, the patch hunk containing the matched comment was inspected, and the smallest justified code unit connected to the comment was selected. When the patch exposed a complete function, method, class, helper, or similar structural unit, it was treated as an entity block. When only a compact changed region could be justified, it was treated as a local code span. File additions and comment-only cases were also recorded when applicable. This process produced 35,278 comments with associated code blocks from 26,490 files in 12,944 repositories.

For RQ1, a stratified sample of 500 comments and code blocks was manually annotated. The annotation process used open coding and axial coding along two dimensions: task type and AI contribution type. Task type describes the development activity supported by AI, such as code implementation, code enhancement, testing, documentation, or bug identification and fixing. AI contribution type describes how AI supported the task, such as implementation, knowledge and concept support, or artifact generation. False positives, such as incidental mentions of names like Claude, were removed. Inter-annotator agreement was measured using Gwet's AC1 \cite{gwet2014handbook}, producing 0.760 for task type and 0.657 for AI contribution type.

To scale annotation to the full dataset, two LLM classifiers, gemma-4:31b and nemotron-3-super:120b, independently labeled the comments and code blocks using few-shot prompts based on the manual taxonomy. Their outputs were consolidated using the Dawid-Skene expectation-maximization framework, which estimates both the most likely true label and the reliability of each annotator \cite{dawid1979maximum}. A set of 400 human-labeled examples initialized the aggregation process, and 100 examples were reserved for held-out evaluation. The held-out agreement was substantial to strong, with AC1 of 0.7005 for task type and 0.8128 for contribution type.

After aggregation, 7,697 instances were identified as false positives and removed. The final labeled corpus contained 27,581 comments with associated code blocks. For task type, the main categories were code implementation, code enhancement, bug identification and fixing, testing, documentation, and generic or indeterminate actions. For AI contribution type, the categories were implementation, knowledge and concept support, artifact generation, and generic or indeterminate actions.

For RQ2, first-change commits were retrieved to analyze how AI-assisted code changed after introduction. For each extracted code block, the corresponding line interval was tracked forward from the introductory commit through later file history. The earliest later commit whose patch overlapped the tracked interval was recorded as the first-change commit. In total, 12,996 first-change commits were retrieved. After removing false positives and very short commit messages, 10,139 first-change commit messages remained for analysis.

BERTopic was used to identify patterns in these first-change commit messages \cite{grootendorst2022bertopic}. Commit messages were preprocessed, embedded using BAAI/bge-base-en-v1.5, reduced using UMAP, and clustered using HDBSCAN. Topic representations were generated with class-based TF-IDF. Because fine-grained topics in short commit messages can be unstable, semantically similar topics were manually grouped into higher-level developer action categories.

For RQ3, monthly time series were constructed using each comment's introductory commit date from December 2022 to March 2026. Category counts were normalized by monthly totals, generic and indeterminate labels were excluded, and the resulting series were smoothed and corrected for anomalies. Descriptive statistics and trend slopes were then computed to examine how task types and AI contribution types evolved over time.

\section{Results and Discussion}

\subsection{RQ1: AI-Assisted Tasks and Contributions}

For RQ1, the results show that AI is primarily used for code implementation. In the full task-specific dataset, code implementation accounts for 18,149 cases, or 72.04\%, as shown in Table~\ref{tab:task_distribution}. This indicates that developers frequently use AI to generate functional code that becomes part of production repositories. Code enhancement is the second-largest task category, with 4,039 cases, or 16.03\%. This shows that AI is also used to improve existing code by increasing readability, changing behavior, improving performance, or refining earlier implementations.

\begin{table}[!t]
\caption{Developer-Level Task Distribution of AI Usage}
\label{tab:task_distribution}
\centering
\small
\begin{tabularx}{0.86\linewidth}{Xrr}
\toprule
\textbf{Task Type} & \textbf{Count} & \textbf{Percent} \\
\midrule
Code Implementation & 18,149 & 72.04\% \\
Code Enhancement & 4,039 & 16.03\% \\
Testing & 1,322 & 5.25\% \\
Documentation & 1,295 & 5.14\% \\
Bug Identification and Fixing & 388 & 1.54\% \\
\midrule
Total & 25,193 & 100.00\% \\
\bottomrule
\end{tabularx}
\end{table}

The AI contribution analysis shows a similar pattern. Implementation is the dominant contribution type, appearing in 16,137 cases, or 82.13\% of the contribution-specific dataset, as shown in Table~\ref{tab:contribution_distribution}. Knowledge and concept support appears in 2,798 cases, or 14.76\%, showing that developers also use AI as an advisor for design decisions, algorithm selection, secure practices, API behavior, syntax, performance issues, and debugging strategies. Artifact generation appears in 714 cases, or 3.63\%, and includes supporting materials such as lists, configuration-like content, documentation text, and test-related artifacts.

\begin{table}[!t]
\caption{Distribution of AI Contribution Types}
\label{tab:contribution_distribution}
\centering
\small
\begin{tabularx}{0.86\linewidth}{Xrr}
\toprule
\textbf{Contribution Type} & \textbf{Count} & \textbf{Percent} \\
\midrule
Implementation & 16,137 & 82.13\% \\
Knowledge and Concept Support & 2,798 & 14.76\% \\
Artifact Generation & 714 & 3.63\% \\
\midrule
Total & 19,649 & 100.00\% \\
\bottomrule
\end{tabularx}
\end{table}

These findings indicate that AI is not limited to a passive recommendation role. In many cases, AI acts as a direct contributor by producing code that developers incorporate into repositories. At the same time, the presence of knowledge and concept support shows that developers also use AI for reasoning and exploration. AI-assisted development therefore includes both execution-oriented and cognition-oriented support.

\subsection{RQ2: Post-Integration Changes}

For RQ2, first-change commits show that AI-assisted code often undergoes substantial post-integration modification. The most frequent action is feature integration and extension, with 2,769 commit messages, as shown in Table~\ref{tab:first_change}. This suggests that AI-assisted code is often incorporated into continuing feature development. Refactoring and cleanup is the second-largest category, with 1,465 messages, including restructuring, formatting, removing unused code, reorganizing imports, and improving design. Bug fixes and corrective changes are also common, with 1,304 messages, indicating that AI-assisted code often requires later correction or adjustment.

\begin{table}[!t]
\caption{Modification Activities in First-Change Commit Messages}
\label{tab:first_change}
\centering
\small
\begin{tabularx}{0.86\linewidth}{Xr}
\toprule
\textbf{Modification Intent} & \textbf{Count} \\
\midrule
Feature Integration and Extension & 2,769 \\
Refactoring and Cleanup & 1,465 \\
Bug Fixes and Corrective Changes & 1,304 \\
Configuration, Dependencies, and Environment & 810 \\
Documentation & 709 \\
Testing and Evaluation & 432 \\
Data, Schema, and Pipeline Processing & 128 \\
Logging and Monitoring & 72 \\
\bottomrule
\end{tabularx}
\end{table}

Other post-integration categories include configuration, dependency, and environment management; documentation; testing and evaluation; data, schema, and pipeline processing; and logging and monitoring. Together, these results show that AI-assisted code is rarely a finished artifact at the moment of introduction. It is usually integrated into the normal software maintenance cycle, where developers extend, refactor, test, document, and fix it.

\subsection{RQ3: Evolution of AI Usage}

For RQ3, AI-referenced activity increases substantially over time. Only 0.29\% of records occur in 2022, followed by 7.73\% in 2023, 21.06\% in 2024, 50.94\% in 2025, and 19.99\% in early 2026. Since the 2026 data covers only through March, the 2026 share represents a partial-year window rather than a full-year decline.

The longitudinal analysis reveals an important shift in AI usage behavior. Code implementation remains the dominant task type, with a mean monthly proportion of 0.720, but it shows a statistically significant decreasing trend over time. Code enhancement increases significantly, suggesting that developers are increasingly using AI to refine and improve existing implementations rather than only generate new code. Testing also increases significantly, indicating growing use of AI in validation and quality-assurance work. Documentation and bug identification do not show statistically significant trends.

\begin{table}[!t]
\caption{Longitudinal Statistics of AI Usage}
\label{tab:longitudinal}
\centering
\small
\begin{tabularx}{\linewidth}{Xrrrrr}
\toprule
\textbf{Category} & \textbf{$\mu$} & \textbf{$\sigma$} & \textbf{$\rho_1$} & \textbf{$\beta$} & \textbf{$p$} \\
\midrule
Code Implementation & 0.720 & 0.065 & 0.884 & -0.00357 & $<$.001 \\
Code Enhancement & 0.191 & 0.039 & 0.799 & 0.00185 & $<$.001 \\
Documentation & 0.043 & 0.026 & 0.497 & 0.00008 & .834 \\
Testing & 0.035 & 0.022 & 0.808 & 0.00119 & $<$.001 \\
Bug ID and Fixing & 0.013 & 0.007 & 0.566 & 0.00012 & .259 \\
Implementation & 0.797 & 0.064 & 0.822 & -0.00324 & .002 \\
Knowledge Support & 0.150 & 0.050 & 0.888 & 0.00286 & $<$.001 \\
Artifact Generation & 0.048 & 0.023 & 0.560 & 0.00052 & .162 \\
\bottomrule
\end{tabularx}
\end{table}

At the contribution level, implementation remains the most common form of AI support, but it also decreases significantly over time. Knowledge and concept support increases significantly, suggesting that developers are gradually using AI more as a reasoning partner and less exclusively as a direct code generator. Artifact generation remains stable and secondary.

Overall, the findings suggest that AI-assisted development is maturing. Early usage is heavily concentrated on direct code generation, while later usage becomes more diverse and includes enhancement, testing, explanation, and conceptual support. This does not mean that implementation becomes unimportant. Instead, it shows that AI's role expands from producing code snippets toward supporting broader development reasoning and maintenance activities.

These results have several implications. First, when AI influences design or implementation decisions, the reasoning behind those decisions should be recorded in project artifacts such as pull requests, issues, or design notes. Otherwise, important knowledge from AI-assisted reasoning may be lost. Second, productivity evaluation should consider the full lifecycle of AI-assisted code, including generation, integration, review, testing, refactoring, and maintenance. Measuring only time to first implementation may overstate the benefit of AI-generated code. Third, AI-assisted development should be understood as a multi-stage practice. Developers may begin by using AI for direct implementation, but over time they also use it for conceptual guidance, improvement, and quality support \cite{rogers2014diffusion}.

\section{Limitations}

This study has several limitations. First, it captures only self-admitted AI usage in code comments. If developers used AI but did not mention it, those cases are not included, so actual AI usage is likely underestimated. Second, keyword-based search can introduce false positives, such as incidental mentions of names or terms unrelated to AI tools. Manual validation and LLM-assisted filtering reduce this risk but cannot eliminate it completely. Third, the dataset focuses only on public GitHub repositories in Python and JavaScript, which may not represent private projects, industrial development, or other programming languages. Fourth, linking comments to code blocks and later commits is approximate because commits can contain unrelated changes, and extracted code blocks may include both AI-assisted and non-AI-assisted code. Finally, AI tools and developer practices are evolving quickly, so the findings represent the observed period from December 2022 to March 2026 rather than a permanent pattern.

\section{Conclusion and Future Work}

This study analyzes AI-referenced GitHub code comments, associated code blocks, and later commits to understand how AI is used in real-world software development. The results show that developers most frequently use AI for code implementation, but also use it for enhancement, testing, documentation, bug fixing, conceptual guidance, and artifact creation. AI-assisted code is commonly followed by feature extension, refactoring, cleanup, bug fixing, documentation, and testing, indicating that human oversight remains essential after AI-supported code enters a project.

The longitudinal results show that AI usage behavior is changing. Direct implementation remains dominant, but knowledge and concept support, code enhancement, and testing become more prominent over time. This suggests that AI-assisted development is not only about generating code. It is increasingly an ongoing human-AI workflow in which developers use AI for reasoning, refinement, and quality support.

Future work should expand the analysis beyond Python and JavaScript and include private or industrial repositories. Additional repository artifacts, such as pull request discussions, issue threads, code review comments, and AI conversation histories, could help reconstruct the full decision-making process behind AI-assisted code. Future studies should also measure the cost and benefit of AI assistance across the full software lifecycle, including generation, integration, review, testing, maintenance, and defect correction. Another important direction is to study whether better documentation of AI-assisted reasoning improves maintainability and knowledge transfer over time.

\bibliographystyle{IEEEtran}
\bibliography{ref}

\end{document}
